% ============================================================
% Section 2: Related Work
% ============================================================

\section{Related Work}
\label{sec:related}

\subsection{Document and Image Forgery Detection}

Early work on image forensics relied on hand-crafted features such as noise
inconsistencies, JPEG compression artifacts, and chromatic aberration patterns to
detect copy-move and splicing manipulations~\cite{huh2018fighting}. Deep learning
substantially shifted the field: convolutional neural networks can learn to detect
subtle statistical fingerprints left by image-processing operations, often
outperforming engineered pipelines on standard benchmarks~\cite{bappy2019hybrid,
mantranet2019}.

A complementary research thread exploits camera-model fingerprints embedded in
image noise. Noiseprint~\cite{cozzolino2020noiseprint} trains a CNN via Siamese
contrastive learning to extract camera-model-specific noise residuals; forgeries
appear as anomalies in this residual and can be localized without semantic
understanding. \textsc{TruFor}~\cite{trufor2023} generalizes this idea with a
transformer fusion network (NoisePrint++) that jointly models RGB semantics and
noise-pattern anomalies, achieving strong performance on general image forensics
benchmarks with over 876\,000 pre-training images. General-purpose detectors such
as MMFusion~\cite{mmfusion2024}, which combines complementary forensic filters
(SRM, Bayar convolution) through an early-and-late fusion decoder, and
FatFormer~\cite{fatformer2024}, which adapts a CLIP-ViT backbone with
forgery-aware spatial and frequency adapters, further advance the state of the art
for natural-image manipulation. UniFD~\cite{unifd2023} demonstrates that
nearest-neighbor classification in a frozen CLIP feature space, without any
discriminative fine-tuning, already achieves competitive detection rates across
diverse generative models.

Identity document forgery, however, introduces unique challenges compared to
natural-image manipulation: (i)~the forged regions are often small (a single
text field or a face) relative to the full document, (ii)~the background texture
is highly structured (guilloché patterns, UV-reactive ink), which creates strong
false negatives for general-purpose noise detectors, and (iii)~the diversity of
document types across countries and issuers demands generalization beyond what
natural-image pre-training provides.
The SIDTD dataset~\cite{sidtd2024} provides a first large-scale benchmark for
this domain and shows that models trained on one set of nationalities degrade
severely on unseen ones, motivating few-shot and domain-generalization approaches.
The FakeIDet line of work~\cite{fakeidet2025} further explores patch-based
DINOv2~\cite{dosovitskiy2021vit} representations with multi-head self-attention fusion for
privacy-preserving fake-ID detection. The DeepID 2025 Challenge~\cite{deepid2025challenge}
was specifically designed to benchmark progress on joint detection and localization
using the FantasyID dataset~\cite{fantasyid2024}, which covers face-swap and
text-inpainting attacks on fantasy identity documents. Official evaluations on
FantasyID show that the general-purpose baselines (TruFor, MMFusion, FatFormer,
UniFD) all achieve false-negative rates above 50\% at a 10\% false-positive rate
threshold~\cite{fantasyid2024}, confirming that document-agnostic forensic models
require domain-specific adaptation.

Among systems trained exclusively on FantasyID, George and
Marcel~\cite{edgedoc2025} propose EdgeDoc, a lightweight hybrid
CNN-Transformer encoder (EdgeNeXt~\cite{edgenext2022}) that takes both the
green channel of the document image and a NoisePrint residual as input, feeding
a U-Net decoder optimized jointly for classification and pixel-level
localization. Naseeb~\textit{et~al.}~\cite{twoheadswin2026} propose
TwoHead-SwinFPN, a Swin Transformer backbone with a feature pyramid network
decoder, CBAM attention modules, and dual-output heads for joint classification
and segmentation, reporting 88.61\% F1 and 57.24\% Dice on FantasyID.
Yu~\textit{et~al.}~\cite{remtkd2025} introduce Re-MTKD, a reinforcement-learning
multi-teacher knowledge distillation framework using a ConvNeXt~v2 encoder
and UPerNet decoder that leverages specialist teachers for copy-move, splicing,
and inpainting; the winning Sunlight team adapted this framework to identity
documents using over 60\,000 additional external training images, achieving the
best results in both tracks of the DeepID 2025 Challenge.

DocVerify differs from all prior document-forensics systems in three respects:
(i)~it systematically searches the task-weight $\lambda_{\text{mask}}$ via
multi-objective HPO rather than fixing it heuristically; (ii)~it selects the
final configuration via a Pareto-optimal criterion that treats detection and
localization as explicitly competing objectives; and (iii)~it reports
uncertainty through nested cross-validation, providing statistically rigorous
performance estimates rather than single-split results.

\subsection{Multi-Task Learning for Detection and Localization}

Joint detection and localization is a canonical instance of multi-task learning
(MTL)~\cite{caruana1997multitask, ruder2017overview}. In MTL, shared representations
learned for one task regularize learning for related tasks, typically yielding better
generalization than independent training when tasks are correlated~\cite{crawshaw2020multi}.

The dominant approach to combining task-specific losses is scalar
scalarization~\cite{marler2004survey}: a single objective $\mathcal{L} = \sum_t
\lambda_t \mathcal{L}_t$ is minimized, with fixed weights $\lambda_t$ set before
training. Kendall \textit{et al.}~\cite{kendall2018multi} proposed learning $\lambda_t$
jointly with the model weights via homoscedastic uncertainty, but this requires an
additional learning-rate schedule and is sensitive to initialization. Liu \textit{et al.}~\cite{liu2019end}
introduced gradient-normalization to dynamically balance task gradients, at the cost
of additional computation per step. All scalar approaches share a common limitation:
the optimal trade-off between tasks is not established before training, yet the
weights implicitly commit to a fixed trade-off.

In the image forensics domain, U-Net~\cite{unet2015} architectures have been widely
adopted for pixel-level localization because their skip connections preserve spatial
resolution while the encoder captures high-level semantic anomaly features. Several
works combine classification and segmentation branches in a single network, but
typically fix the task weights empirically without systematic
optimization~\cite{wu2022robust}. Our work differs by including the task weight
$\lambda_{\text{mask}}$ in the HPO search space and selecting it via a
multi-objective criterion rather than a grid search.

\subsection{Hyperparameter Optimization and Multi-Objective Search}

Hyperparameter optimization (HPO) has evolved from manual search and grid search to
sequential model-based methods. Tree-structured Parzen Estimator
(TPE)~\cite{tpe2011} models the distribution of good and bad configurations
separately and selects candidates that maximize the ratio of the two models, a
strategy significantly more sample-efficient than random search~\cite{randomsearch2012}.
The Optuna framework~\cite{optuna2019} implements TPE as its default sampler and
has become a standard tool for HPO in deep learning.

Multi-objective HPO extends single-objective search to the case where multiple
metrics must be jointly optimized. The output is a Pareto-optimal set of
configurations rather than a single best point, making the trade-off between
objectives explicit~\cite{miettinen1999nonlinear}. Evolutionary algorithms such as
NSGA-II~\cite{nsga2_2002} are the most established approach, but they require
large population sizes to maintain Pareto-front diversity. The Multiobjective TPE
(MOTPE) sampler~\cite{motpe2022} adapts the TPE surrogate model to the
multi-objective setting, achieving competitive Pareto-front coverage with far fewer
function evaluations than evolutionary methods—a critical advantage when each
evaluation requires training a deep neural network.

To our knowledge, this is the first work to apply Pareto-based multi-objective HPO
to the problem of joint identity-document forgery detection and localization, and
the first to use the resulting Pareto front as the primary model-selection criterion
in a nested cross-validation framework.
